A bőrrák, és különösen a melanoma korai felismerés esetén jól kezelhető, késői
stádiumban azonban romlik a túlélési esély; a bőrgyógyászati
szakellátáshoz való hozzáférés ugyanakkor világszerte egyenetlen.
Éppen ezért jelen dolgozat célja egy mély tanuláson alapuló mesterséges intelligencia modell,
mely képes a bőrelváltozások kép alapú osztályozására, és az e köré épülő webböngészőben és
mobileszközön elérhető, ,,bőrdaganat-szűrő'' rendszer bemutatása, amely a bőrléziók szakorvosi
vizsgálatának indokoltságát jelzi.

Osztályozóként transzfertanulással finomhangolt Vision Transformer (ViT) modellt
használtunk, melyet az ISIC 2024-es adathalmazán tanítottunk, így alkalmassá vált
nem-dermatoszkópos bőrlézió-felvételek klasszifikációjára. A modellt azonos feltételek
mellett betanított ResNet-18 konvolúciós hálóval, az ISIC versenygyőztes EVA02 modelljével
és Gemini 2.5 Flash nagy nyelvi modellel hasonlítottuk össze. 
A finomhangolt ViT a teszthalmazon AUC $0{,}939$-et ért el, érdemben felülmúlva
a konvolúciós alapot, és elérve az EVA02 szintjét.
A betanított osztályozó modellt egy teljes alkalmazássá egészítettük ki: mobil- és webklienssel, a
modellt kiszolgáló MCP-réteggel, a Gemini modellre épülő képvalidációval és
természetes nyelvű magyarázattal, valamint hitelesített, felhasználónként
elkülönített felhőalapú adattárolással.

\textbf{Kulcsszavak}: bőrrák, melanoma, mély tanulás, Vision Transformer,
transzfertanulás, ISIC 2024, nagy nyelvi modell, mobil
egészségügy
